% ege_ref.tex
% Справочные данные ЕГЭ по физике 2027 (из спецификации и демоверсии ФИПИ)
\begingroup
\footnotesize
\setlength{\parskip}{2pt}
\setlength{\tabcolsep}{3.5pt}
\renewcommand{\arraystretch}{0.92}

\section*{Справочные материалы}

\subsection*{Десятичные приставки}
\noindent\begin{tabularx}{\linewidth}{|l|c|c||l|c|c|}
\hline
\textbf{Приставка} & \textbf{Обозначение} & \textbf{Множитель} & \textbf{Приставка} & \textbf{Обозначение} & \textbf{Множитель} \\ \hline
гига  & Г & $10^9$  & санти & с & $10^{-2}$ \\ \hline
мега  & М & $10^6$  & милли & м & $10^{-3}$ \\ \hline
кило  & к & $10^3$  & микро & мк & $10^{-6}$ \\ \hline
гекто & г & $10^2$  & нано  & н & $10^{-9}$ \\ \hline
деци  & д & $10^{-1}$ & пико & п & $10^{-12}$ \\ \hline
\end{tabularx}

\subsection*{Константы и нормальные условия}
\noindent\begin{tabularx}{\linewidth}{>{\raggedright\arraybackslash}X r}
Ускорение свободного падения на Земле & $g = 10~\text{м/с}^2$ \\
Гравитационная постоянная & $G = 6{,}67 \cdot 10^{-11}~\text{Н}\cdot\text{м}^2/\text{кг}^2$ \\
Скорость света в вакууме & $c = 3{,}0 \cdot 10^8~\text{м/с}$ \\
Элементарный заряд & $e = 1{,}6 \cdot 10^{-19}~\text{Кл}$ \\
Постоянная Планка & $h = 6{,}63 \cdot 10^{-34}~\text{Дж}\cdot\text{с}$ \\
Постоянная Авогадро & $N_A = 6{,}02 \cdot 10^{23}~\text{моль}^{-1}$ \\
Универсальная газовая постоянная & $R = 8{,}31~\text{Дж/(моль}\cdot\text{К)}$ \\
Постоянная Больцмана & $k = 1{,}38 \cdot 10^{-23}~\text{Дж/К}$ \\
Нормальное атмосферное давление & $p_0 = 10^5~\text{Па} = 760~\text{мм рт. ст.}$ \\
Нормальная температура & $T_0 = 273~\text{К}~(0~^\circ\text{C})$ \\
\end{tabularx}

\subsection*{Соотношения между единицами и масса частиц}
\noindent\begin{tabularx}{\linewidth}{X r || X r}
\hline
1 электрон-вольт (эВ) & $1{,}6 \cdot 10^{-19}~\text{Дж}$ & Масса протона $m_p$ & $1{,}673 \cdot 10^{-27}~\text{кг}$ \\ \hline
1 атомная единица массы (а.е.м.) & $1{,}66 \cdot 10^{-27}~\text{кг}$ & Масса нейтрона $m_n$ & $1{,}675 \cdot 10^{-27}~\text{кг}$ \\ \hline
1 киловатт-час (кВт$\cdot$ч) & $3{,}6 \cdot 10^6~\text{Дж}$ & Масса электрона $m_e$ & $9{,}1 \cdot 10^{-31}~\text{кг}$ \\ \hline
\end{tabularx}

\subsection*{Плотности веществ}
\noindent\begin{tabularx}{\linewidth}{X r || X r}
\hline
Вода & $1000~\text{кг/м}^3$ & Керосин & $800~\text{кг/м}^3$ \\ \hline
Лёд & $900~\text{кг/м}^3$   & Ртуть & $13600~\text{кг/м}^3$ \\ \hline
Воздух (при н.у.) & $1{,}29~\text{кг/м}^3$ & Медь & $8900~\text{кг/м}^3$ \\ \hline
Сталь, железо & $7800~\text{кг/м}^3$ & Алюминий & $2700~\text{кг/м}^3$ \\ \hline
\end{tabularx}

\subsection*{Молярные массы веществ}
\noindent\begin{tabularx}{\linewidth}{X r || X r}
\hline
Водород ($\text{H}_2$) & $2 \cdot 10^{-3}~\text{кг/моль}$ & Кислород ($\text{O}_2$) & $32 \cdot 10^{-3}~\text{кг/моль}$ \\ \hline
Гелий ($\text{He}$) & $4 \cdot 10^{-3}~\text{кг/моль}$ & Воздух & $29 \cdot 10^{-3}~\text{кг/моль}$ \\ \hline
Азот ($\text{N}_2$) & $28 \cdot 10^{-3}~\text{кг/моль}$ & Вода ($\text{H}_2\text{O}$) & $18 \cdot 10^{-3}~\text{кг/моль}$ \\ \hline
\end{tabularx}

\subsection*{Удельная теплоёмкость}
\noindent\begin{tabularx}{\linewidth}{X r || X r}
\hline
Вода & $4200~\text{Дж/(кг}\cdot\text{К)}$ & Лёд & $2100~\text{Дж/(кг}\cdot\text{К)}$ \\ \hline
Алюминий & $900~\text{Дж/(кг}\cdot\text{К)}$ & Медь & $400~\text{Дж/(кг}\cdot\text{К)}$ \\ \hline
Железо & $460~\text{Дж/(кг}\cdot\text{К)}$ & Свинец & $130~\text{Дж/(кг}\cdot\text{К)}$ \\ \hline
\end{tabularx}

\subsection*{Удельная теплота плавления и парообразования}
\noindent\begin{tabularx}{\linewidth}{X r || X r}
\hline
Плавление льда & $3{,}3 \cdot 10^5~\text{Дж/кг}$ & Парообразование воды & $2{,}3 \cdot 10^6~\text{Дж/кг}$ \\ \hline
Плавление свинца & $2{,}5 \cdot 10^4~\text{Дж/кг}$ & Сгорание бензина & $4{,}6 \cdot 10^7~\text{Дж/кг}$ \\ \hline
\end{tabularx}
\vspace{6pt}
\endgroup
\newpage
